\section{Discussion}
\label{sec:discussion}

\paragraph{Without instructions.} What is forced is the conflict of the inputs of two tasks
(Proposition~\ref{prop:noinstructionprice}); the rest is not. A recogniser on the frozen model's own
states avoids most of the loss, blocks no task and does not depend on the order of arrival; on fresh tasks it passed the line fixed for it, as did two related rules on other fresh tasks. Two tasks
built on the same texts split their inputs under it, and \AudConfPairShare{}\% of their inputs occur
word for word in both, where no method without instructions can do better. With the published number
of epochs the two methods learn the five tasks nearly as well as plain adapters with the same number of epochs (\AudFullInfloraLearned{}
and \AudFullGainloraLearned{} right after learning, against \AudFullOursRidgeLearned{}) and then forget; the
forgetting grew with what had been learned. Plain adapters behind the recogniser keep the same routes
and the same \AudFullOursRidgeChanged{} of \AudFullCompared{} changed answers at 10 and at 100 epochs. With their definitions the same run gives \AudFullDefOriginalInfloraFT{} (InfLoRA) and \AudFullDefOriginalGainloraFT{} (GainLoRA) points of forgetting, with \AudFullDefOriginalInfloraChanged{} and \AudFullDefOriginalGainloraChanged{} of \AudFullDefCompared{} answers changed: in this setting the instruction is what held most of their answers (nearly all for GainLoRA, fewer for InfLoRA), and our rules do not need it.
On fifteen tasks without definitions (one seed, our rules only) they end
\AudFifteenNoDefRidgeBelowOwn{} points below the ceiling with \AudFifteenNoDefRidgeChanged{} of
\AudFifteenNoDefCompared{} changed answers.

\paragraph{Two rules, two kinds of statement.} The zones give an exact statement inside them, none
outside, and read the instruction; the recogniser gives a measured statement for all inputs of a task
and, with \handicap{}, a loose one in advance, and reads only the example. With instructions the two
behaved as stated: the zones changed no answer and lost \AudUThreeHeavyCertBelowOwn{} points on average
under the heavy rewording through misrouted inputs, and the recogniser stayed within
\AudUThreeRidgeBelowOwnHigh{} points of the ceiling (Section~\ref{sec:exampleroutes}). Using the zone
where an input is certified and the recogniser elsewhere gave the better of the two at all wordings but
the medium one (seed 42; Table~\ref{tab:instanceroutes}). A system should use the zones if it needs a
proof for unedited instructions, and the recogniser if its routes must not depend on the wording.

\paragraph{The third rule.} Theorem~\ref{thm:bound} needs no instruction: with each zone calibrated on
held-out inputs of its task, an input changes its answer with probability at most a chosen level, for
any model and any later tasks. The first tests (Section~\ref{sec:noinstruction}: routes on benchmark files without definitions, and a small trained model) met the bound and showed its
price: with frozen word-vector features the zones of earlier tasks hold most inputs of a new task, which
is then blocked and not learned. The difficulty moves to one quantity measured at admission, the share
of a task's inputs inside older zones. Telling tasks apart without an instruction is what the
recogniser takes up (Sections~\ref{sec:audit} and~\ref{sec:recogniser}).

\paragraph{Comparison with the published methods, with instructions.} It would be easy to read the gap
between the certified switch and InfLoRA and GainLoRA as ``they forget and we do not''. On five tasks
that is wrong: both have forgetting scores near zero (though \FiveInfChanged{} and \FiveGainChanged{} of
\FiveInfCompared{} answers change), and they had learned less, at the level of guessing on the two tasks
with two labels. On fifteen tasks InfLoRA forgets \FifteenInfFT{} points and GainLoRA \FifteenGainFT{},
and both had also learned less (\FifteenInfLearned{} and \FifteenGainLearned{} against
\FifteenCertLearned{}), so there the gap has a part from learning and a part from forgetting. We did not find out why: restricting an adapter to a subspace, or training only one of its two matrices, may slow learning,
ten epochs may be too few, and our set-up replaced one step of their code. Plain adapters learn these tasks well and then lose them; the switch removes the loss without touching the learning. What the experiments
support is narrow: switching protects without changing how an adapter is trained, at the price of
collisions with old zones and, at scale, of worse routing from the fallback. They do not support a statement about how
well the published methods retain what they learn: the five-task runs left little to retain, the 15-task run has one seed, the run with definitions at the full number of epochs has one seed (Section~\ref{sec:nodeftrained}), and none used a larger model. Without their definitions the two methods did hold their answers less well (Section~\ref{sec:nodefpublished}), with the same limits.

\paragraph{What the zones add to a nearest-task rule.} On the benchmark sequences \nearest{} and
\cert{} give the same routes and scores, because no later task tries to take an earlier task's inputs.
They separate in two situations. When a later task resembles an old one and inputs of the old task
arrive in other words, best match can hand them to the newcomer without notice and the zones keep those
inside a zone (Section~\ref{sec:stress}, a constructed case). When all instructions are unedited, best
match serves both tasks and the zones refuse the newcomer: in \ManyTasks{} real tasks this happened to
\ManyCollidingPct{}\% (Section~\ref{sec:many}). On machine paraphrases of \WidenTasksEval{} real tasks
this case was rare and the two rules gave the same routes (Section~\ref{sec:paraphrase}); the zones then
added the proof and not better routing, and a plain nearest-signature rule gave more paraphrased inputs to their own task. We therefore cannot say that a system that expects reworded instructions is better served by the zones, only that it gets a proof for part of them. Sibling tasks with fixed instructions are better served by best
match, by a smaller radius, or by the variant of Section~\ref{sec:exactfirst}, which gives up the guarantee only for inputs that carry the sibling's instruction word for word. What the zones offer in
both situations is that a conflict is known when the task arrives, as the count of blocked training
inputs, and that a certified answer is final.

\paragraph{Exactness cuts both ways.} A certified answer never gets worse and never better. InfLoRA and
GainLoRA show small negative forgetting on some tasks (an old sentiment task improves when a second one
is learned); the switch cannot profit. An improved version of an old task learned as a new task would
meet the old opening of its inputs, so the old adapter keeps answering them under every zone rule here
and, under \recog{}, the two versions would split them (Proposition~\ref{prop:noinstructionprice}).
Replacing a claim, with an explicit decision to give up the guarantee for one zone, is the natural
correction; we have not studied it.

\paragraph{Averaged word-vectors.} Two instructions that share few content words are far apart, and an
edit that keeps most content words stays close, so light rewordings are certified. A heavy rewording
can leave few content words, and then tasks about the same subject are hard to tell apart: the only
inputs the zones route wrongly are tweets under a one-sentence instruction about mood, given to the
other sentiment task. Features that read meaning (hidden states, a sentence encoder, an inference model
as judge; Section~\ref{sec:paraphrase}, Appendix~\ref{app:widen}) raise the share of correct
certificates, but none brings the wrong certificates under the 5\% asked for: the lowest,
\WidenEncRulerTwentyFrM{}\% with \WidenEncRulerTwentyFrP{}\% correct (\WidenRefTfiveFrM{}\% and
\WidenRefTfiveFrP{}\% for \cert{}), belongs to a setting singled out after the runs and beat \cert{}
by less than the margin fixed on held-out paraphrases. Sibling tasks that differ in one decisive word
are put together as well. The theorem needs only zones fixed when the task arrives and membership
computed reproducibly by frozen parts (Remark~\ref{rem:otherzones}).

\paragraph{Mechanical edits.} The edit test (Section~\ref{sec:edits}) shows three things. Robustness
depends on length: an instruction of a hundred tokens keeps its route under sixteen typos, one of ten
words loses it under eight, so tasks should have instructions of a few sentences. The switch reads a
bag of words: shuffled sentences are certified, and a removed sentence matters only if it carried the
words that set the task apart. Typos are worse than missing words, because a misspelt word is split
into rare word pieces that the weighting of Eq.~\eqref{eq:weights} takes for content; a
spelling-tolerant feature could be swapped in without touching the guarantee, but we did not test one.

\paragraph{Uncertified inputs.} For a certified input the answer is final. For an uncertified one the
system knows it is guessing the task, and $\min_k \delta_k(x)$, the distance to the nearest zone in
units of its radius, indicates how good the guess is (we did not test how well it predicts a correct
route). The system can answer anyway, ask for the instruction in its registered form, or refuse; we did
not study these policies, which matter for badly damaged instructions, where more than half of the
routes are wrong and all are flagged.

\paragraph{The benchmark hands every method the task.} Every input of SuperNI begins with its task
definition, so the task can be read off by string matching, and methods described as needing no task
identity receive it in the text. The exact key reaches the full learning of plain adapters with no
forgetting and collapses to plain adapters when one word of the instruction changes. Evaluations of
continual instruction tuning should include reworded instructions, or a lookup table is a perfect
method (Sections~\ref{sec:paraphrase} and~\ref{sec:nodefpublished}).

\paragraph{Evaluation noise.} The official loader scores an input against one reference picked at
random on every load, so identical answers differ by about a point on a task with several references.
On the 15 tasks this gave the certified switch a forgetting of \FifteenCertOffFT{} although none of
\FifteenCertCompared{} answers changed (Section~\ref{sec:fifteen}). When the claim is that nothing
changed, answers should be compared, not scores.

\section{Limitations}
\label{sec:limits}

\begin{itemize}
  \item \textbf{Reduced setting.} T5-base instead of T5-large and 10 epochs instead of 100 in all runs
  but a few, on free GPUs. The two published methods were at the level of guessing on two of the five
  tasks right after learning them and learned less than plain adapters on the 15 tasks, and one step of
  their code was a substitute for its library, without a number-for-number check
  (Appendix~\ref{app:impl}). Nothing follows about these methods in the published setting.
  \item \textbf{Comparison without instructions.} Two published methods and one sequence of five tasks:
  three seeds at a tenth of the epochs (most of the gap is learning) and one seed at the full number,
  where our rules end at \AudFullOursRidgeAPOne{} and \AudFullOursHandicapAPOne{}. The rules were
  chosen on the 15 benchmark tasks, so those numbers, and the five GPU tasks, are development data.
  Numbers on fresh tasks, on the second benchmark and with 120 tasks are routes only. Trained answers without instructions exist
  for our rules only on fifteen tasks (one seed), and not on the second benchmark or with the published
  methods given an oracle or our recogniser (Section~\ref{sec:nodeftrained}).
  \item \textbf{The bound in advance is loose:} the handicap gives \AudDevHPrimeTTwoBound{}\% on average
  against \AudDevHPrimeTTwoChanged{}\% measured. The split of fresh tasks does not ensure that test and
  held-out inputs are the same kind of draw (a random split moved the numbers by
  \AudScaleSplitMaxDiffEnd{} points or less).
  \item \textbf{Shared inputs and people.} Without instructions no method can do better than split
  inputs that occur in two tasks (\AudScalePairOlderLEnd{}\% and \AudScalePairNewerLEnd{}\% of the
  inputs of two tasks on the same texts reached their own task). Rewordings by people were not tested.
  \item \textbf{Routes from the example alone:} five tasks, three seeds, machine-written development
  rewordings (Section~\ref{sec:exampleroutes}). The benchmark gives instruction and example as two
  fields; finding where the example begins in one string was not tested.
  \item \textbf{Chosen tasks.} The five-task sequence was put together from tasks that one published
  method had forgotten most in an earlier run of ours, and the unrelated inputs come from a benchmark
  without its sentiment tasks. The first choice was not varied, the second only for routes
  (Section~\ref{sec:background}).
  \item \textbf{Latent-LoRA only as our re-implementation.} It has no public code; we ran it with a third
  of its epochs, it learned little, so we compare no scores; its router gave every input to the right task and no answer changed. It was not run in the stress test
  (Appendix~\ref{app:latent}).
  \item \textbf{Few seeds, small test sets.} Three seeds for five tasks, one for fifteen; about 100
  inputs per task, so differences of a few points on one task are not meaningful, and two runs with the
  same seed differ by a few tenths of a point.
  \item \textbf{At scale the switch adds the proof, not better routing.} On machine paraphrases of
  \WidenTasksEval{} real tasks the nearest-zone rule changed \RerouteNearNChanged{} of
  \RerouteInputs{} routes after the tasks arrived and the certified switch the same
  \RerouteCertNChanged{}, against the \RerouteExpectNearLow{}\% to \RerouteExpectNearHigh{}\% we had
  expected. With the fallback of all GPU runs the switch routed worse than a plain nearest-signature rule
  (\RerouteCertOwnEnd{}\% against \RerouteSigOwnEnd{}\% to their own task). A variant of the fallback,
  scaled by each task's distance to unrelated text and tested at the end on routes only, is right on the
  benchmark (\RerouteBgHeavyFive{}\% under heavy rewording) and at scale (\RerouteBgItOwn{}\% without
  zones and \RerouteCertFourItOwn{}\% with zones in front, through a fourth language), but the best rule
  without zones still gives \RerouteZonesCost{} to \RerouteZonesCostEs{} points more right routes, depending on the pivot language. The two fallbacks trade places between
  tests, and the variant has not been run with trained answers (Section~\ref{sec:paraphrase}).
  \item \textbf{The zones are narrow at scale.} Under machine paraphrases \cert{} certifies \WidenRefTfivePLow{}\% to \WidenRefTfivePHigh{}\% of the
  inputs for their own task and \WidenRefTfiveMLow{}\% to \WidenRefTfiveMHigh{}\% for another, depending
  on the pivot language; about three quarters reach their own task, against \FiveCertRightLight{}\%,
  \FiveCertRightMedium{}\% and \FiveCertRightHeavy{}\% on the five benchmark tasks, which are far apart.
  None of five attempts to widen the zones (hidden states, a paraphrase ruler, a sentence encoder with a
  conformal radius, an inference model as judge, the same judge sentence by sentence) met its bar fixed
  before the runs: correct certificates rose, but wrong ones rose with them (hidden states, paraphrase
  ruler) or stayed above the bar (sentence encoder \WidenEncPrimM{}\%, judge \WidenJudgeFrM{}\%,
  sentence-by-sentence judge \WidenSentItM{}\% with \WidenSentItP{}\% correct on the Italian pivot,
  against 9\% for a failure and 5\% for a success; Appendix~\ref{app:widen}). A sentence-encoder zone
  chosen after the runs was ahead of \cert{} on each measure on held-out paraphrases, but by less than
  the margin fixed in advance and not on every measure in other splits, so we do not claim it.
  \item \textbf{Development rewordings.} The reworded instructions were written by machine during
  development, share words with the originals, and two design choices (the radius factor, and the rule
  for inputs outside the zones with the word weights) were made while looking at where they were routed.
  An independent test needs rewordings by people who have not seen the method.
  \item \textbf{The guarantee covers certified inputs only.} With medium and heavy rewording most inputs
  are outside the zones; they kept their answers in our runs, but a later look-alike task can take them
  (stress test), and how many reach their own task depends on the unrelated inputs that set the scale:
  \BackgroundHeavyLow{}\% to \BackgroundHeavyHigh{}\% under heavy rewording (Section~\ref{sec:background}).
  Under heavy damage, sixteen typos leave only 45\% of the inputs on their own task and very short
  instructions lose their route after a few edits; these inputs are flagged, not rescued.
  \item \textbf{The switch ignores word order.} Its feature is a weighted average of the first tokens, so two instructions of the same words in another order cannot be told apart.
  \item \textbf{Several tests are about routes only:} machine paraphrases, re-routing, the five attempts to
  widen the zones, the Spanish check, the stress test, mechanical edits, task orders and other embedding
  tables. They do not measure what a model answers once an input has reached its adapters, and the
  paraphrase tests use one split and one arrival order (except the five splits of the Spanish check).
  \item \textbf{Zones from the inputs of a task block later tasks.} The bound of Theorem~\ref{thm:bound}
  holds for test inputs like the training inputs of their task and was met in every run
  (\BoundAllChangedLow{}\% to \BoundAllChangedHigh{}\% of routes changed against the \BoundAlpha{}\%
  allowed), but most later tasks were blocked on arrival (Section~\ref{sec:noinstruction}).
  \item \textbf{The instruction zones need an instruction.} They are built from the common opening of a
  task's inputs; a task without one has no signature and is served by \recog{}
  (Sections~\ref{sec:zones}, \ref{sec:recogniser} and~\ref{sec:nodeftrained}).
  \item \textbf{Older tasks have priority.} A new task that collides with an old zone is blocked on the
  colliding inputs; among \ManyTasks{} real tasks almost all collisions were siblings with a nearly
  identical instruction, and \ManyAllHeld{} tasks could not be learned at all. The remedies are
  rewording the new instruction, a smaller radius factor or giving up the old claim; \exactfirst{}
  serves such tasks with a weaker guarantee and was tested on routes only. On unedited instructions a
  nearest-zone rule gives more inputs to their own task (\ManyNearRightPct{}\% against
  \ManyCertRightPct{}\%). Certified answers cannot improve through later tasks (no backward transfer).
  \item \textbf{One trained model family, and memory.} Training and answers were tested with T5-base and,
  on a CPU, with the stand-alone code on small models; other models only for routes. One adapter per task
  is kept, as in all methods of this family.
  \item \textbf{The scores do not separate the switch from simpler rules:} on the benchmark sequences a
  nearest-zone rule, and with unedited instructions an exact match of the instruction, reach the same
  scores; the certified switch adds the proof and the certificate, not points.
\end{itemize}

\section{Conclusion}
\label{sec:conclusion}

With one frozen set of adapters per task, an answer changes only if its route changes, so forgetting is
a question about the rule that sets the route. Without instructions every method loses at least the
conflict of two tasks, and benchmarks hide this because their inputs begin with the instruction. The
published methods we ran lost answers when it was removed, and this was not forced. Our recogniser blocks
no task, is the same for every order of arrival and comes with an audit and a loose bound in advance;
plain adapters behind it beat the two methods without their instructions in a reduced setting by a line fixed before the run, with
the gap being learning at a tenth of their epochs and forgetting at the published number (one seed); the five tasks are among the development tasks. Its routes are
the same for every wording of the instruction, where the zones lose inputs under heavy rewording but change no answer.

When inputs carry an instruction, the oldest claim on a zone wins and the answer to an input inside it
cannot change, for any model and any later tasks. In a reduced setting on a real model, on five tasks and on 15, this kept every answer while learning
as much as unprotected adapters, with the original instructions and, in the order we ran, with reworded ones, but simpler rules reach the same scores; the zones add the proof. Their
limits: the guarantee ends at the border of the zones, they are built from words and not meaning, and a
new task that collides with an old claim is blocked. Missing are the larger model of the published
setting, more seeds at the full number of epochs, the published methods on more than five tasks without instructions, a tight
bound in advance, rewordings by people, and tasks that share inputs.
